\thispagestyle{empty}\begingroup\fontsize{9}{11}\selectfont\raggedright\noindent{\Large\bfseries Translation sources and terminology\par}\medskip\noindent The Latin text is from Carl Friedrich Gauss, \textit{Werke}, Band IV (Göttingen, 1873), issued by the Royal Society of Sciences at Göttingen. The original-language editable transcription and its historical portrait are available in the \href{https://github.com/KokunoYumeto/gauss-collected-works/releases/tag/gauss-original-language-2026-10-24}{original-language edition}. The modern source encoding and typesetting are credited to OpenAI GPT-6 Astra, Ultra; this English translation and its typesetting are AI-produced.\par\medskip\noindent Related English mathematical-argument usage was consulted in James Caddall Morehead and Adam Miller Hiltebeitel’s \textit{General Investigations of Curved Surfaces of 1827 and 1825} (1902), \href{https://www.gutenberg.org/ebooks/36856}{Project Gutenberg 36856, TeX edition}. Related German quantity and mathematical-argument usage was consulted in H. Maser’s translation and edition, \textit{Carl Friedrich Gauss’ Untersuchungen über höhere Arithmetik} (1889), \href{https://archive.org/details/carlfriedrichga00gausgoog}{historical digitized edition}. These works are not parallel translations of the observation theory. Gauss’s German notice of 26 February 1821, in the same historical volume, was consulted for related error and probability terminology. No human running prose or formulas are adopted from these references. Historical publication titles remain in the original. Exact versions and URLs are retained in \texttt{CITATIONS.bib}.\par\endgroup\clearpage
